% 示例阅读顺序：查看实际效果，在源码中复制对应 frame，再按需替换内容。
\section{引文与文献}
\subsection{引用文字与出处}
\begin{frame}{引文与引用}
  \begin{quote}
    \enquote{引文内容写在这里，通常是一句原文；\texttt{quote} 环境整段缩进，
    适合较短的引文。}
  \end{quote}

  \vspace{0.3em}
  \begin{itemize}
    \item 叙述式引用：\textcite{example2026} 提出了该问题的基本框架。
    \item 括号式引用：已有工作给出了类似结论\parencite{example2026}。
    \item 只取作者与年份：\citeauthor{example2026}（\citeyear{example2026}）。
    \item 长引文用 \texttt{\textbackslash blockquote}：\blockquote{这是一段
      较长的引文，超过一定行数后会自动转为整段缩进，并在末尾给出出处。}
  \end{itemize}
  \sourceNote{来源：\parencite{example2026}（示例条目，请替换为真实文献）。}
\end{frame}

\begin{frame}{多段引文与普通脚注}
  \begin{quotation}
    第一段引文：quotation 适合包含多个段落的引用文字。

    第二段引文：空行分段，保留引文环境的缩进和专用字体。
  \end{quotation}
  正文中的补充说明可以使用普通脚注\footnote{这是一条补充说明；脚注自动编号。}。
  \sourceNote{出处使用 sourceNote；需要文献时可加入 \parencite{example2026}。}
\end{frame}

\subsection{脚注引用与文献表}
% 普通脚注用 footnote，整条文献用 footcite；文献键来自 references.bib。
\begin{frame}{脚注式引用}
  \begin{itemize}
    \item 作者 + 年份：此处引用一条工作\footnote{\citeauthor{example2026},
      \citeyear{example2026}.}。
    \item 作者 + 标题 + 年份：此处引用另一条工作\footnote{\citeauthor{example2026},
      \citetitle{example2026}, \citeyear{example2026}.}。
    \item 完整文献条目：用 \texttt{\textbackslash footcite} 直接把文献表内容
      放进脚注\footcite{example2026}。
  \end{itemize}

  \begin{alertblock}{使用建议}
    \begin{itemize}
      \item 图表与数据来源用 \texttt{\textbackslash sourceNote}；
      \item 观点出处用 \texttt{\textbackslash parencite} / \texttt{\textbackslash textcite}；
      \item 只在需要展开出处细节（如封面引文）时使用 \texttt{\textbackslash footnote}。
    \end{itemize}
  \end{alertblock}
\end{frame}

\begin{frame}[fragile]{文献表与附录的写法}
  \begin{verbatim}
\appendix
\begin{frame}[allowframebreaks]{参考文献}
  \printbibliography[heading=none]
\end{frame} % 结束参考文献页面
  \end{verbatim}
  \slideNote{在 references/references.bib 中添加条目，正文通过键名引用。
  附录页码独立编号；实际的参考文献表在文档末尾展示。}
\end{frame}
